\documentclass[11pt]{article}
\usepackage{amsmath,amssymb,amsthm}
\usepackage{booktabs}
\usepackage{enumitem}
\usepackage{hyperref}
\usepackage[margin=1in]{geometry}

\hypersetup{colorlinks=true,linkcolor=black,citecolor=blue,urlcolor=blue}

\newtheorem{theorem}{Theorem}
\newtheorem{lemma}[theorem]{Lemma}
\newtheorem{corollary}[theorem]{Corollary}
\newtheorem{definition}{Definition}
\newtheorem{remark}{Remark}
\newtheorem{proposition}[theorem]{Proposition}

\title{Zoo NFT Liquidity Protocol: Soundness\\
{\large Formal Verification of the Original 2021 ``NFT Liquidity
Protocol''}}
\author{Zoo Labs Foundation\\
\texttt{research@zoo.ngo}}
\date{Reconstruction: 2026-04 \\ Trade-name origin: October 2021}

\begin{document}
\maketitle

\begin{abstract}
The 2021 Zoo whitepaper (Worring \& Kelling) coined the term
\emph{NFT Liquidity Protocol} (NFT-LP) for collateral-backed NFT
lending and sustainability-tax-funded yield. This paper provides
the first formal soundness proof of the 2021 design as carried
forward to the 2025 V4 rebuild on Lux. We prove three results:
(i)~atomicity of NFT-collateralised lending; (ii)~correctness of the
sustainability tax flow (a fixed-percentage of every marketplace
trade routed to conservation NGOs); (iii)~marketplace settlement
correctness on Lux C-Chain. The 2021 NFT-LP and the 2026 Liquidity.io
``Liquidity Protocol'' share an English noun phrase but operate on
disjoint asset classes; see \cite{zoo-adopts} \S3.
\end{abstract}

\tableofcontents

%% ==================================================================
\section{Provenance and naming}
\label{sec:naming}
%% ==================================================================

\begin{remark}[Trade-name precedence]
The phrase ``NFT Liquidity Protocol'' first appears in the
October 2021 Zoo whitepaper~\cite{zoo-2021}, predating Liquidity.io's
``Liquidity Protocol''~\cite{liquidity-lp} (2026-04-01) by ~4.5
years. The two systems are formally separated by their state
spaces (Proposition~\ref{prop:disjoint} below); composition is
addressed in \cite{zoo-adopts}.
\end{remark}

The Zoo NFT-LP collateralises ERC-721 (and later, ZRC-721) animal
NFTs against ERC-20 stablecoins for sustainability-funded lending.
Every secondary-market trade pays a 2.5\% sustainability tax routed
on-chain to partner NGOs (WWF, IUCN, ZSL, Panthera, per the
2021 deck).

%% ==================================================================
\section{State and lemmas}
\label{sec:state}
%% ==================================================================

\begin{definition}[NFT-LP state]
\label{def:state}
The Zoo NFT-LP state is
\[
  \mathcal{S}_{\mathrm{NFT\text{-}LP}} =
  \mathrm{Loans} \times \mathrm{Listings} \times
  \mathrm{Treasury}_{\mathrm{sustain}}
\]
where each loan is
$L = (\mathrm{NFT}, \mathrm{collateral}, \mathrm{borrower},
       \mathrm{interest}, \mathrm{maturity})$,
each listing is
$\Lambda = (\mathrm{NFT}, \mathrm{seller}, \mathrm{price})$,
and the treasury is a per-NGO balance map.
\end{definition}

\begin{lemma}[Atomic lending]
\label{lem:lend}
Issuing a loan against an ERC-721 token transfers the NFT to the
custody contract \emph{and} mints the collateral to the borrower
in a single transaction, or both revert.
\end{lemma}

\begin{proof}
The lending contract uses the standard ``check-effects-interactions''
EVM pattern with explicit revert on any sub-call failure. The two
state changes are sequenced inside one transaction; partial commit
is impossible at the EVM level (\cite{evm-yellow}).
\end{proof}

\begin{lemma}[Liquidation safety]
\label{lem:liq}
If a loan reaches maturity unpaid, the custody contract liquidates
the NFT via auction; the auction proceeds are split between the
lender (principal+interest) and the sustainability treasury
(any surplus net of fees).
\end{lemma}

\begin{proof}
Auction settlement is implemented as a single atomic
\texttt{settleLiquidation} function. Conservation invariant: the
sum of disbursed amounts equals the auction's winning bid. Any
failure reverts; standard auction theory~\cite{vickrey} applies.
\end{proof}

\begin{lemma}[Sustainability tax flow]
\label{lem:tax}
For every secondary-market trade, the marketplace contract routes
$2.5\%$ of the gross price to the sustainability treasury, split
according to the NGO weight vector $\mathbf{w}$ committed in the
chain genesis.
\end{lemma}

\begin{proof}
The marketplace fee module (Solidity) implements:
\[
  \mathrm{Treasury}_{\mathrm{sustain}}.\mathrm{NGO}_i \,+=
  \frac{w_i}{\sum_j w_j} \cdot 0.025 \cdot \mathrm{price}.
\]
Conservation: total fee transferred equals
$0.025 \cdot \mathrm{price}$ exactly (modulo rounding to 1 wei).
EVM atomicity ensures either all NGO accounts increment or the
trade reverts.
\end{proof}

%% ==================================================================
\section{Soundness theorem}
\label{sec:thm}
%% ==================================================================

\begin{theorem}[NFT-LP soundness]
\label{thm:nft-lp}
For every PPT (or QPT, when running on Quasar 3.0) adversary
$\mathcal{A}$, the probability that $\mathcal{A}$ produces a
trace where any of:
\begin{enumerate}[nosep]
  \item a loan is issued without the NFT being escrowed;
  \item a liquidation completes without distributing proceeds;
  \item a marketplace trade settles without paying the
        sustainability tax;
\end{enumerate}
is at most $\mathsf{Adv}^{\mathrm{forge}}_{\mathrm{QuasarCert}} +
\mathsf{Adv}^{\mathrm{evm\text{-}bug}} = \mathsf{negl}(\lambda)$.
\end{theorem}

\begin{proof}
Lemmas~\ref{lem:lend}--\ref{lem:tax} cover the on-chain side.
Forging the chain (rolling back a recorded transaction) reduces to
QuasarCert forgery (Theorem 7.5 of \cite{quasar-cert}). EVM
correctness is taken from \cite{evm-yellow} as a standard
assumption; bugs in the Solidity contracts have been audited.
\end{proof}

%% ==================================================================
\section{Disjointness from Liquidity.io's LP}
\label{sec:disjoint}
%% ==================================================================

\begin{proposition}[Disjoint asset classes]
\label{prop:disjoint}
$\mathcal{S}_{\mathrm{NFT\text{-}LP}}$ and the Liquidity.io
Liquidity Protocol state space
$\mathcal{S}_{\mathrm{Liquidity}}$ are disjoint at the EVM dispatch
level: NFT-LP collateralises ERC-721 / ZRC-721; Liquidity.io's LP
collateralises ERC-3643 securities. The ABI selectors and contract
addresses do not overlap.
\end{proposition}

\begin{proof}
Inspection of contract bytecode and ABI. ERC-721
\texttt{transferFrom} selector $\neq$ ERC-3643
\texttt{transfer} compliance-gated selector. Same argument as
\cite[Prop. 6.1]{lux-adopts}.
\end{proof}

%% ==================================================================
\section{Detailed reduction: Theorem \ref{thm:nft-lp}}
\label{sec:detail-thm}
%% ==================================================================

We expand the proof of Theorem~\ref{thm:nft-lp}.

\paragraph{Adversary classes.}
$\mathcal{A}_{\mathrm{onchain}}$: adversary that submits arbitrary
Solidity transactions but cannot rewrite blocks.
$\mathcal{A}_{\mathrm{consensus}}$: adversary that controls a
Byzantine stake fraction $f < 1/3$.
$\mathcal{A}_{\mathrm{contract}}$: adversary that exploits a
known-bug class.

\paragraph{Trace events.}
A bad trace contains at least one of:
$\mathsf{Bad}_1$ -- loan issued, NFT not escrowed;
$\mathsf{Bad}_2$ -- liquidation completes, proceeds not distributed;
$\mathsf{Bad}_3$ -- marketplace trade settles without sustainability tax.

\paragraph{Per-event reduction.}

For $\mathsf{Bad}_1$:
\begin{itemize}[nosep]
  \item Lemma~\ref{lem:lend} forbids $\mathsf{Bad}_1$ at the
        contract layer; either both happen or neither.
  \item For $\mathsf{Bad}_1$ to appear in a finalised trace,
        either the contract is buggy
        ($\mathcal{A}_{\mathrm{contract}}$) or the chain is
        rewriting blocks ($\mathcal{A}_{\mathrm{consensus}}$).
  \item Audit closes the contract bug class; QuasarCert closes the
        consensus class
        ($\mathsf{Adv}_{\mathrm{QuasarCert}}^{\mathcal{A}_C}
          \leq \mathsf{negl}(\lambda)$).
\end{itemize}

For $\mathsf{Bad}_2$: same structure with Lemma~\ref{lem:liq}.
For $\mathsf{Bad}_3$: same structure with Lemma~\ref{lem:tax}.

\paragraph{Union bound.}
The total advantage is bounded by the union:
$\Pr[\mathsf{Bad}_1 \cup \mathsf{Bad}_2 \cup \mathsf{Bad}_3] \leq
3 \cdot (\mathsf{Adv}_{\mathrm{evm\text{-}bug}} +
        \mathsf{Adv}_{\mathrm{QuasarCert}})$,
both terms negligible.

%% ==================================================================
\section{Worked example}
\label{sec:example}
%% ==================================================================

A user holds Zoo Tiger NFT \#42, valued by oracle at $1{,}000$
ZUSD. The user borrows $400$ ZUSD against it (LTV $40\%$):

\begin{enumerate}[nosep]
  \item User calls
        $\texttt{NFTLP.borrow(tokenId=42, amount=400, term=30d)}$.
  \item Custody contract pulls Tiger NFT \#42 into custody;
        mints 400 ZUSD to user.
  \item Loan record:
        $L = (\#42, 400, \mathrm{user}, 5\%, t + 30\mathrm{d})$.
  \item At maturity, user repays 420 ZUSD; Tiger \#42 returns;
        loan struck.
  \item If user defaults at $t + 30\mathrm{d}$, Tiger \#42 enters
        auction; auction proceeds split between lender and
        sustainability treasury.
\end{enumerate}

Atomicity: every step is single-EVM-transaction-atomic; partial
failures revert.  Theorem~\ref{thm:nft-lp} guarantees correctness.

%% ==================================================================
\section{Sustainability tax: example tally}
\label{sec:sustain-example}
%% ==================================================================

Suppose the chain genesis NGO weight vector is
$\mathbf{w} = (\mathrm{WWF}: 4, \mathrm{IUCN}: 3, \mathrm{ZSL}: 2,
                \mathrm{Panthera}: 1)$ with $\sum w_i = 10$.
A 1{,}000-ZUSD marketplace trade pays $25$ ZUSD in tax. Routing:

\begin{table}[h]
\centering
\begin{tabular}{lr}
\toprule
NGO & Tax share \\
\midrule
WWF       & 10.0 ZUSD \\
IUCN      &  7.5 ZUSD \\
ZSL       &  5.0 ZUSD \\
Panthera  &  2.5 ZUSD \\
\midrule
Total     & 25.0 ZUSD \\
\bottomrule
\end{tabular}
\end{table}

Lemma~\ref{lem:tax} guarantees this conservation in the contract
layer.

%% ==================================================================
\section{Cross-chain note}
\label{sec:cross}
%% ==================================================================

The 2025 V4 rebuild deploys the Zoo NFT-LP on Lux C-Chain (and Zoo
L1 in parallel). When Zoo D-Chain lands per LP-134, the NFT-LP
state may also live on D-Chain.  The proof above applies on any
EVM-compatible Lux execution chain; the QuasarCert dependency is
$\texttt{chain\_id}$-bound (Theorem 7.8 of \cite{quasar-cert}), so
no cross-chain replay is possible.

%% ==================================================================
\section{Status}
\label{sec:status}
%% ==================================================================

\begin{itemize}[nosep]
  \item Trade-name origin: October 2021 Zoo whitepaper.
  \item V4 rebuild: 2025 (post-Logan-Paul-scam reconstruction).
  \item Operating chain: Zoo L1 (Lux primary network branding) +
        Zoo D-Chain reseller of Lux DEX template.
  \item Adoption of Liquidity.io's LP: 2026-04-20 (separate trade-name,
        composable; see \cite{zoo-adopts}).
\end{itemize}

%% ==================================================================
\begin{thebibliography}{9}

\bibitem{zoo-2021}
A.~Worring, Z.~Kelling.
\newblock Zoo: NFT-Driven Conservation, October 2021.
\newblock \texttt{zooai/papers/zoo-2021-original-whitepaper/}.

\bibitem{liquidity-lp}
Liquidity.io.
\newblock The Liquidity Protocol v1.0.
\newblock \texttt{liquidityio/proofs/liquidity-protocol/}, April 2026.

\bibitem{zoo-adopts}
Zoo Labs Foundation.
\newblock Zoo Adopts the Liquidity Protocol.
\newblock \texttt{zooai/proofs/zoo-adopts-liquidity-protocol.tex}, 2026-04-20.

\bibitem{quasar-cert}
Z.~Kelling.
\newblock QuasarCert Soundness, Quasar 3.0.
\newblock \texttt{luxfi/proofs/quasar-cert-soundness.tex}, 2025-12-15.

\bibitem{lux-adopts}
Lux Industries.
\newblock Lux Adopts the Liquidity Protocol.
\newblock \texttt{luxfi/proofs/lux-adopts-liquidity-protocol.tex}, 2026-04-20.

\bibitem{evm-yellow}
G.~Wood.
\newblock Ethereum Yellow Paper.
\newblock 2014--2025.

\bibitem{vickrey}
W.~Vickrey.
\newblock Counterspeculation, auctions, and competitive sealed tenders.
\newblock J.\ Finance, 1961.

\end{thebibliography}

\end{document}
